% !TEX encoding = UTF-8
% !TEX TS-program = pdflatex
% !TEX root = ../tesi.tex

\chapter{Conclusione}
\label{cap:cap4}

\intro{In questo capitolo viene attuata una valutazione retrospettiva del tirocinio svolto, andando a stilare gli obiettivi raggiunti e le conoscenze acquisite. Viene infine ritagliato uno spazio per le considerazioni personali sull'esperienza del tirocinio.}\\

\section{Obiettivi raggiunti}
Al completamento del tirocinio, insieme al \textit{project manager} e al \textit{tutor} aziendale, abbiamo verificato il soddisfacimento degli obiettivi stabiliti. Quasi tutti tutti gli obiettivi sono stati raggiunti, in particolare tutti gli obiettivi obbligatori sono stati rispettati ed è stato soddisfatto anche l'unico obiettivo desiderabile, andando a implementare funzioni inizialmente non previste, come il calcolo delle distanze tra i partecipanti, il \textit{caching} dei dati e la visualizzazione della \textit{route} percorsa. La tabella \ref{tab:consu} offre un riepilogo dello stato degli obiettivi alla conclusione del tirocinio.

\begin{longtable}{ |p{2cm}|p{6.5cm}|p{2.5cm}|  }
        \caption{Completamento degli obiettivi del tirocinio.}
        \label{tab:consu}
        \hline
        \textbf{Codice} & \textbf{Descrizione} & \textbf{Stato}
        \hline
        \textbf{OO1} & La gestione e la pianificazione del progetto deve essere gestita tramite una \textit{board kanban} condivisa. & raggiunto \\ \hline
        \textbf{OO2} & Il tirocinante deve praticare un'analisi e una progettazione dell'applicazione \textit{mobile} a partire dai requisiti raccolti.& raggiunto \\ \hline
        \textbf{OO3} & Il tirocinante deve sviluppare il modulo di autenticazione dell'applicazione.& raggiunto \\ \hline
        \textbf{OO4} & Il tirocinante deve sviluppare il modulo di raccolta e invio dei dati \textit{gps}. & raggiunto\\ \hline
        \textbf{OO5} & Il tirocinante deve sviluppare il modulo per la gestione dei \textit{tour}, \textit{route} e \textit{POI}.& raggiunto\\ \hline
        \textbf{OD1} & Il tirocinante deve essere in grado di coordinarsi con il cliente.& raggiunto \\ \hline
        \textbf{OD2} & Il tirocinante deve sviluppare la \textit{suit} di \textit{testing} del prodotto.& non raggiunto \\ \hline
        \textbf{OD3} & Il tirocinante deve redarre una documentazione completa del prodotto.& raggiunto \\ \hline
        \textbf{OP1} & Ulteriori modifiche al prodotto non comprese negli obiettivi obbligatori. & raggiunto\\ \hline
\end{longtable}

\section{Conoscenze acquisite}
Il periodo di tirocinio mi ha permesso di apprendere nuovi processi di sviluppo, a me prima sconosciuti. Partendo dall'inizio dell'esperienza: durante le prime settimane ho acquisito una buona dimetichezza con il linguaggio \textit{Dart} e il \textit{framework Flutter}, contemporaneamente, sono stato immerso nella metodologia di lavoro agile, di cui non conoscevo le modalità prima di allora.\\
Nelle settimane successive ho raffinato le conoscenze tecniche andando nel dettaglio di nuovi \textit{pattern} architetturali e qui ho appreso per la prima volta il \textit{pattern bloc}, che nel progetto ho poi usato estensivamente. Lavorare su un'applicazione multipiattaforma mi ha poi introdotto nello specifico agli ambienti \textit{iOS} e \textit{Android}, con la quale mi sono interfacciato per il processo di rilascio dell'applicativo.\\
Infine, dal progetto in sè ho assimilato concetti specifici, come l'utilizzo di \textit{API} di mappe geografiche, interfacciamento ai sistemi \textit{GPS} e l'utilizzo di formati \textit{GPX}.

\section{Considerazioni}
\subsection{L'azienda}
L'esperienza con \textit{Moku S.r.l.}, l'azienda ospitante, è stata sicuramente positiva Durante i due mesi di tirocinio il \textit{team} è rimasto sempre dispobibile a qualsiasi chiarimento e dubbio. L'unica difficoltà che posso dire di aver incontrato è stato lo studio di \textit{flutter}, infatti la tecnologia in questione era nuova per il \textit{tean}, quindi l'apprendimento iniziale è stato a mio parere più lungo di quanto sarebbe potuto essere nel caso in cui ci fosse stato un esperto a cui rivolgermi.

\subsection{Il cliente}
L'esperienza con il cliente è stata una successione di alti e bassi, sebbene il cliente fosse sempre disponibile a eventuali chiarimenti e discussioni, \textit{CyclingHero} è stato un definito troppo velocemente e questo a causato diversi rallentamenti: in primo luogo i requisiti sono stati più volte rivisitati dal cliente, creando inutili perdite di tempo, in secondo la tempistica per i rilasci è stata troppo frettolosa, incentivando uno sviluppo poco solido, volto solamente a rispettare le scadenze. Infine, al termine del tirocinio il \textit{back-end} offerto dal cliente si trovava ancora in una stadio preliminare, ciò mi ha impedito di implementare diverse funzionalità, in particolare quella di autenticazione, propedeutica a molte altre funzionalità.

\subsection{Commento finale}
L'esperienza del tirocinio mi portato a una crescita su più livelli, che oltre alle conoscenze tecniche acquisite mi ha permesso di inquadrare per la prima volta un'ambiente di lavoro reale, attraverso interazioni tra colleghi, \textit{project manager} e clienti esterni.